\documentclass[../../../include/open-logic-section]{subfiles}

\begin{document}

\olfileid{sth}{z}{story}
\olsection{વધુ વિગતવાર વર્ણન}

\olref[story][approach]{sec}માં આપણે ગણરચનાની પ્રક્રિયા અંગેનું
શોએનફેલ્ડનું વર્ણન અવતરિત કર્યું હતું. હવે આ વર્ણનને થોડું વધુ
ચોક્કસ બનાવવા માટે થોડા વધુ સિદ્ધાંતો લખવા ઇચ્છીએ છીએ. તે આ છે:
\begin{enumerate}
\item[] \stageshier. દરેક ગણ કોઈક તબક્કે રચાય છે.
\item[] \stagesord. તબક્કાઓ ક્રમબદ્ધ છે: કેટલાક બીજા તબક્કાઓની
\emph{પહેલાં} આવે છે.\footnote{ખરેખર આપણે અપ્રત્યક્ષ રીતે
માનીશું કે તબક્કાઓ \emph{સુક્રમિત} છે. તેનો અર્થ શું થાય તે
\olref[ordinals][]{chap}માં સમજાવ્યું છે. આ એક મહત્ત્વની ધારણા છે.
હકીકતમાં \citet{Scott1974}ની ખૂબ ચતુર પદ્ધતિ વાપરીને આ ધારણાને
\emph{ટાળી} શકાય અને પછી તેને \emph{તારવી} શકાય. (આ પ્રથમ તબક્કો
હોવો જોઈએ એમ શા માટે માનવું તે પણ સમજાવશે.) અહીં તેની વિગતમાં
જઈ શકતા નથી; વધુ માટે \citet{ButtonLT1} જુઓ.}
\item[] \stagesacc. કોઈ પણ તબક્કા $S$ માટે, અને તબક્કા~$S$ની
\emph{પહેલાં} રચાયેલા કોઈ પણ ગણો માટે: તબક્કા~$S$ પર એવો ગણ
રચાય છે જેના સભ્યો બરાબર તે ગણો જ હોય. તબક્કા~$S$ પર બીજું
કંઈ રચાતું નથી.
\end{enumerate}
આ અનૌપચારિક સિદ્ધાંતો છે, પરંતુ તેમની મદદથી ઝર્મેલોના
ગણસિદ્ધાંતની કેટલીક સ્વયંસિદ્ધિઓને યોગ્ય ઠરાવી શકીશું.

(એક સાવચેતી જણાવવી જોઈએ. આપણે સંપૂર્ણપણે માનક નામોવાળી
સંપૂર્ણપણે માનક સ્વયંસિદ્ધિઓ રજૂ કરીશું, પરંતુ હમણાં રજૂ કરેલા
ત્રાંસા અક્ષરોવાળા સિદ્ધાંતો માટે સાહિત્યમાં કોઈ ખાસ નામો નથી.
આપણે માત્ર એવા નામો આપ્યાં છે જે મદદરૂપ થશે એવી આશા છે.)

\end{document}
